\documentclass[11pt,a4paper]{article}
\usepackage[utf8]{inputenc}
\usepackage[T1]{fontenc}
\usepackage{amsmath,amssymb,amsfonts}
\usepackage{amsthm}
\usepackage[margin=2.5cm]{geometry}
\usepackage{enumitem}
\usepackage{tikz}
\usepackage{xcolor}
\usepackage{hyperref}
\hypersetup{colorlinks=true,linkcolor=blue!55!black,urlcolor=blue!55!black}
\setlist[itemize]{leftmargin=1.4em,itemsep=2pt,topsep=3pt}

% ── theorem environments (numbered, standard) ──
\newtheorem{theorem}{Theorem}[section]
\newtheorem{lemma}[theorem]{Lemma}
\newtheorem{proposition}[theorem]{Proposition}
\newtheorem{corollary}[theorem]{Corollary}
\theoremstyle{definition}
\newtheorem{definition}[theorem]{Definition}
\newtheorem{example}[theorem]{Example}
\newtheorem{remark}[theorem]{Remark}
\newtheorem{exercise}{Exercise}

% ── master macro block (safe, multi-character) ──
\newcommand{\Z}{\mathbb{Z}}
\newcommand{\Q}{\mathbb{Q}}
\newcommand{\R}{\mathbb{R}}
\newcommand{\N}{\mathbb{N}}
\newcommand{\Pow}{\mathcal{P}}
\newcommand{\id}{\operatorname{id}}

\title{
  {\color{blue!40!black}\textbf{\Huge Point}}\\[0.3em]
  \Large \textit{From Number Theory to Algebraic Geometry and Logic}\\[0.7em]
  \large \textbf{Session P0 --- Set Theory, Order, and the Axiom of Choice}\\[0.3em]
  \normalsize \emph{Complete lecture notes with proofs, examples, and exercises}
}
\author{}
\date{}

\begin{document}
\maketitle
\thispagestyle{empty}

\begin{abstract}
\noindent
This is the true session zero of the course --- the one the first class actually needed. It builds
the set-theoretic scaffolding everything else stands on: equivalence relations and their exact
correspondence with partitions (the logic of every quotient to come: $\Z/n$, $G/N$, $M\otimes N$,
sheafification); order relations (preorder, partial order, strict order) and the anatomy of
extrema (minimum vs minimal, bounds, well-orders); and the axiomatic frame ZF with the
\textbf{Axiom of Choice}, its equivalents (Well-Ordering, Zorn's lemma --- with the equivalence
cycle proved or sketched), and its famous consequences and equivalents: bases of vector spaces,
Tychonoff, Tarski's theorem, spanning trees, and the precise tier of the Completeness Theorem.
Zorn's lemma is the engine of three existence theorems later in the course: Krull (S1), algebraic
closures (S3), and Deligne (S58).
\end{abstract}

\subsection*{Session plan (150 minutes)}
\begin{center}\small
\begin{tabular}{|l|l|}
\hline
\textbf{Time} & \textbf{Block} \\
\hline
0:00--0:10 & Orientation: why set theory first; where P0 feeds (S1, S3, S2, S27, S58) \\
0:10--0:35 & \S1 Equivalence relations and partitions (the logic of quotients) \\
0:35--1:00 & \S2--\S3 Orders: preorder/partial/strict; minima vs minimal; bounds; well-orders \\
1:00--1:10 & \emph{Break} \\
1:10--1:35 & \S4 From Russell's paradox to ZF \\
1:35--2:10 & \S5 Axiom of Choice; the equivalence cycle; the list of equivalents \\
2:10--2:30 & Exercises / wrap-up: Zorn as the course's existence engine \\
\hline
\end{tabular}
\end{center}

\subsection*{Conventions}
A \emph{relation} on $X$ is a subset of $X\times X$; we write $x\sim y$ for $(x,y)\in R$.
Orders are written $\preceq$ (reflexive form) and $\prec$ (strict form). ``Poset'' means
partially ordered set. ZF denotes Zermelo--Fraenkel set theory without Choice; ZFC $=$ ZF $+$ AC.

\tableofcontents
\newpage

% ══════════════════════════════════════════════
\section{Equivalence Relations and Partitions}
% ══════════════════════════════════════════════

\begin{definition}[Equivalence relation; class; quotient]
A relation $\sim$ on $X$ is an \textbf{equivalence relation} if it is reflexive ($x\sim x$),
symmetric ($x\sim y\Rightarrow y\sim x$), and transitive. The \textbf{class} of $x$ is
$[x]=\{y\in X:y\sim x\}$; the \textbf{quotient set} is $X/{\sim}=\{[x]:x\in X\}$.
\end{definition}

\begin{definition}[Partition]
A \textbf{partition} of $X$ is a family $\mathcal P\subseteq\Pow(X)$ of nonempty, pairwise disjoint
sets whose union is $X$.
\end{definition}

\begin{theorem}[Equivalence $\leftrightarrow$ partition]\label{thm:equiv-partition}
\begin{enumerate}[label=(\alph*)]
  \item If $\sim$ is an equivalence relation on $X$, then $X/{\sim}$ is a partition of $X$.
  \item If $\mathcal P$ is a partition of $X$, then $x\sim_{\mathcal P} y\iff \exists P\in\mathcal P,\ x,y\in P$
    is an equivalence relation, and its classes are exactly the members of $\mathcal P$.
\end{enumerate}
Hence equivalence relations on $X$ and partitions of $X$ are in canonical bijection.
\end{theorem}
\begin{proof}
(a) Each $x\in[x]$ by reflexivity, so the union is $X$ and classes are nonempty. If
$z\in[x]\cap[y]$ then $x\sim z$ and $z\sim y$, so $x\sim y$; then for any $u\in[x]$: $u\sim x\sim y$
gives $u\in[y]$, hence $[x]\subseteq[y]$; symmetrically $[y]\subseteq[x]$. So distinct classes are
disjoint. (b) Reflexive since $\mathcal P$ covers $X$; symmetric by definition; transitive because
if $x,y\in P$ and $y,z\in Q$ then $y\in P\cap Q$, forcing $P=Q$ by disjointness, so $x,z\in P$.
The class of $x$ is the unique $P\ni x$. The two passages invert each other.
\end{proof}

\begin{example}[Quotients everywhere]\label{ex:quotients}
\begin{itemize}
  \item Congruence mod $n$ on $\Z$: classes $[0],\dots,[n-1]$; quotient $\Z/n\Z$ (S1).
  \item Construction of $\Q$: on $\Z\times(\Z\setminus\{0\})$, $(a,b)\sim(c,d)\iff ad=bc$;
    checking this is an equivalence relation \emph{is} the construction of the rationals.
  \item Cosets $G/H$ (P1), orbits of a group action (P1), stalks and sheafification (S26--S27),
    tensor products as quotients of free modules (S2): every quotient in the course is
    Theorem~\ref{thm:equiv-partition} at work --- \emph{identifying what should be the same point}.
\end{itemize}
\end{example}

\begin{remark}[Choosing representatives]
A quotient map $X\to X/{\sim}$ generally has no canonical section: picking one representative per
class for infinitely many classes is exactly where the Axiom of Choice (\S5) enters. Keep this in
mind for S2 and S27.
\end{remark}

% ══════════════════════════════════════════════
\section{Order Relations and Their Kin}
% ══════════════════════════════════════════════

\begin{definition}[Preorder, partial order, strict order, total order]
A relation $\preceq$ on $P$ is a \textbf{preorder} if reflexive and transitive; a \textbf{partial
order} if also antisymmetric ($x\preceq y\preceq x\Rightarrow x=y$); \textbf{total/linear} if any two
elements are comparable. A relation $\prec$ is a \textbf{strict partial order} if irreflexive and
transitive; a \textbf{strict total order} if additionally $x\prec y$ or $x=y$ or $y\prec x$ for all $x,y$.
\end{definition}

\begin{proposition}[The two forms of order correspond]\label{prop:order-forms}
If $\preceq$ is a partial order then $x\prec y\iff(x\preceq y\ \wedge\ x\neq y)$ is a strict partial
order; if $\prec$ is a strict partial order then $x\preceq y\iff(x\prec y\ \vee\ x=y)$ is a partial
order. These passages are mutually inverse.
\end{proposition}
\begin{proof}
Irreflexivity of $\prec$ is immediate; transitivity: if $x\prec y\prec z$ then $x\preceq z$, and
$x=z$ would give $x\preceq y\preceq x$, hence $x=y$ by antisymmetry, contradicting $x\prec y$; so
$x\neq z$, i.e.\ $x\prec z$. The converse direction and the round-trips are similar case analyses.
\end{proof}

\begin{proposition}[A preorder is a partial order on its quotient]\label{prop:preorder-quotient}
If $\preceq$ is a preorder on $P$, then $x\sim y\iff(x\preceq y\ \wedge\ y\preceq x)$ is an
\emph{equivalence relation}, and $[x]\preceq[y]\iff x\preceq y$ is a well-defined \emph{partial order}
on $P/{\sim}$. Thus \S1 and \S2 meet: every preorder becomes a poset after quotienting.
\end{proposition}
\begin{proof}
Reflexivity/symmetry of $\sim$ are built in; transitivity uses transitivity of $\preceq$ twice.
Well-definedness: if $x\sim x'$ and $x\preceq y$ then $x'\preceq x\preceq y\preceq y'$, so
$x'\preceq y'$. Antisymmetry on classes is exactly the definition of $\sim$.
\end{proof}

\begin{example}\label{ex:orders}
\begin{itemize}
  \item $(\N,\mid)$ divisibility: partial, not total ($2\nmid 3$, $3\nmid 2$).
  \item $(\Pow(X),\subseteq)$: partial; chains $=$ nested families, antichains $=$ pairwise disjoint (or incomparable) families.
  \item Preorder of functions $f\preceq g\iff f=O(g)$ (analysis): not antisymmetric --- quotienting gives the poset of growth rates (Proposition~\ref{prop:preorder-quotient}).
  \item Lexicographic order on $\R^2$: total.
\end{itemize}
\end{example}

% ══════════════════════════════════════════════
\section{Extrema, Bounds, and Well-Orders}
% ══════════════════════════════════════════════

\begin{definition}[Bounds and extrema]
In a poset $(P,\preceq)$ with $S\subseteq P$: $u$ is an \textbf{upper bound} of $S$ if $s\preceq u$
for all $s\in S$ (lower bound dually); the \textbf{supremum} is the least upper bound.
$m\in S$ is the \textbf{minimum} of $S$ if $m\preceq s$ for all $s\in S$; $m$ is \textbf{minimal} if
no $s\in S$ satisfies $s\prec m$. Maximum/maximal are dual.
\end{definition}

\begin{example}[Minimum $\neq$ minimal]\label{ex:minimal}
In $P=\{a,b,c\}$ with $a\prec c$, $b\prec c$ and $a,b$ incomparable (Figure~1): the minimal elements
are $a$ and $b$ (two of them!), there is \emph{no} minimum; $c$ is both maximum and maximal; the set
$\{a,b\}$ has upper bound $c$ (its supremum) but no maximum. In $(\N_{\ge2},\mid)$: minimal elements
are the primes; the minimum is $2$? no --- $2$ does not divide $3$: there is no minimum, only minima.
\end{example}

\begin{figure}[h]
\centering
\begin{tikzpicture}[scale=1.2]
  \node (c) at (0,1.6) {$c$};
  \node (a) at (-1,0) {$a$};
  \node (b) at (1,0) {$b$};
  \draw (a) -- (c); \draw (b) -- (c);
  \node at (0,-0.8) {\small two minimal elements, no minimum; $c$ = maximum = supremum of $\{a,b\}$};
\end{tikzpicture}
\caption{Hasse diagram: minimal vs minimum, maximal vs maximum, bounds vs extrema.}
\end{figure}

\begin{definition}[Well-order]
A total order $\preceq$ on $W$ is a \textbf{well-order} if every nonempty $S\subseteq W$ has a
minimum. Examples: $(\N,\le)$; $(\Z,\le)$ and $(\R,\le)$ are \emph{not}; every ordinal is.
\end{definition}

\begin{proposition}\label{prop:wo-basic}
A well-order has no infinite strictly decreasing chain $x_1\succ x_2\succ\cdots$; conversely, a total
order with no infinite decreasing chain is a well-order \emph{provided} countable choice is available
(over ZF the forward direction is unconditional).
\end{proposition}
\begin{proof}
If $x_1\succ x_2\succ\cdots$ exists then $\{x_i:i\in\N\}$ is nonempty with no minimum. Conversely, if
$S\neq\varnothing$ has no minimum, pick $x_1\in S$; having $x_n$, the set $\{s\in S:s\prec x_n\}$ is
nonempty (else $x_n$ were minimal$=$minimum), so choose $x_{n+1}$ in it --- producing a decreasing
chain. (The recursive choices use dependent choice; note where.)
\end{proof}

% ══════════════════════════════════════════════
\section{From Russell's Paradox to ZF}
% ══════════════════════════════════════════════

\begin{proposition}[Russell]
There is no set $R$ with $R=\{x:x\notin x\}$.
\end{proposition}
\begin{proof}
If $R$ were a set, then $R\in R\iff R\notin R$ by definition --- contradiction. Hence
``the set of all sets'' cannot exist: comprehension must be restricted.
\end{proof}

\begin{definition}[The ZF axioms, informally]
Language: only $\in$ and $=$. Axioms: \textbf{Extensionality} (same elements $\Rightarrow$ same set);
\textbf{Pairing} $\{x,y\}$; \textbf{Union} $\bigcup\mathcal F$; \textbf{Power set} $\Pow(X)$;
\textbf{Infinity} (an inductive set exists); \textbf{Separation} (from a set, carve out by any
formula --- the safe replacement of naive comprehension); \textbf{Replacement} (images of sets under
definable functions are sets); \textbf{Foundation} (no infinite $\in$-chains; every nonempty set has
an $\in$-minimal element). \textbf{Choice} is \emph{not} among them --- it is the guest of \S5.
\end{definition}

\begin{remark}[What ZF cannot do alone]
ZF proves each $X_i\neq\varnothing$ has elements, but not that an \emph{infinite family} admits a
simultaneous choice: Russell's socks (no distinguishing property) vs shoes (left/right gives a rule).
Everything in \S5 measures exactly this gap.
\end{remark}

% ══════════════════════════════════════════════
\section{The Axiom of Choice and Its Equivalents}
% ══════════════════════════════════════════════

\begin{definition}[AC, WO, Zorn]
\textbf{AC:} every family $\{X_i\}_{i\in I}$ of nonempty sets has a choice function
$c:I\to\bigcup_i X_i$ with $c(i)\in X_i$; equivalently $\prod_i X_i\neq\varnothing$.
\textbf{WO (Well-Ordering Principle):} every set admits a well-order.
\textbf{Zorn (ZL):} if $(P,\preceq)\neq\varnothing$ and every chain has an upper bound in $P$, then
$P$ has a maximal element.
\end{definition}

\begin{proposition}[WO $\Rightarrow$ AC, in ZF]\label{prop:wo-ac}
\end{proposition}
\begin{proof}
Well-order $W=\bigcup_i X_i$. For each $i$, $X_i$ is a nonempty subset of $W$, hence has a
$W$-least element; set $c(i)=\min_W X_i$. This is a choice function --- no choice used in the proof.
\end{proof}

\begin{proposition}[Zorn $\Rightarrow$ AC, in ZF]\label{prop:zorn-ac}
\end{proposition}
\begin{proof}
Let $\mathcal F$ be the set of partial choice functions: $f\in\mathcal F$ iff $f$ is a function with
$\mathrm{dom}(f)\subseteq I$ and $f(i)\in X_i$. Order $\mathcal F$ by extension; $\mathcal F\neq\varnothing$
(empty function). A chain $\mathcal C$ has upper bound $\bigcup\mathcal C$ (union of compatible
functions is a partial choice function). Zorn gives maximal $f$. If $\mathrm{dom}(f)\neq I$, pick
$i\notin\mathrm{dom}(f)$ and $x\in X_i$ (a \emph{single} choice, legal in ZF) and extend $f$ by
$i\mapsto x$ --- contradicting maximality. Hence $\mathrm{dom}(f)=I$.
\end{proof}

\begin{proposition}[AC $\Rightarrow$ WO and AC $\Rightarrow$ Zorn: sketches]\label{prop:ac-wo-zorn}
\emph{AC $\Rightarrow$ WO:} fix a choice function $c$ on $\Pow(X)\setminus\{\varnothing\}$; by
transfinite recursion set $x_\alpha=c\big(X\setminus\{x_\beta:\beta<\alpha\}\big)$; the recursion must
stop (no injection from the class of ordinals into a set --- Hartogs), and at the stopping stage $X$
is exhausted, so $\alpha\mapsto x_\alpha$ well-orders $X$. \emph{AC $\Rightarrow$ Zorn:} assuming no
maximal element, recursively build a strictly increasing chain $x_\alpha$ by choosing at each stage
a strict upper bound of $\{x_\beta:\beta<\alpha\}$ (which exists since chains have upper bounds and
no upper bound is maximal); again Hartogs bounds the length --- contradiction.
\end{proposition}

\begin{corollary}
Over ZF, the three statements AC, WO, ZL are equivalent (Propositions~\ref{prop:wo-ac},
\ref{prop:zorn-ac}, \ref{prop:ac-wo-zorn}); so are Hausdorff's Maximality Principle and Tukey's lemma.
\end{corollary}

\begin{theorem}[The list of equivalents and near-equivalents]\label{thm:ac-list}
Over ZF, the following are \textbf{equivalent} to AC:
\begin{enumerate}[label=(\arabic*)]
  \item Every vector space has a basis (Blass, 1984).
  \item Tychonoff's theorem: arbitrary products of compact spaces are compact.
  \item Tarski's theorem: for every infinite $A$, $|A\times A|=|A|$.
  \item Every connected graph has a spanning tree (forward direction: Zorn on forests ordered by inclusion).
  \item Every surjection has a section (a reformulation of AC itself).
\end{enumerate}
And the following live in strictly weaker tiers (important to know, easy to overstate):
the \textbf{Boolean Prime Ideal / Ultrafilter Lemma} and the general \textbf{G\"odel Completeness
Theorem} (equivalent to BPI; the countable-language case is provable in ZF); \textbf{Hahn--Banach};
and mere \emph{consequences} of AC such as the existence of non-measurable sets.
\end{theorem}

\begin{remark}[Where the course spends Zorn]
Krull's theorem (S1: maximal ideals exist) $\to$ algebraic closures (S3) $\to$ Deligne's theorem
(S58: coherent topoi have enough points). Three existence theorems, one engine --- and each is a
``points exist'' statement, fitting the course's theme perfectly.
\end{remark}

% ══════════════════════════════════════════════
\section{Exercises}
% ══════════════════════════════════════════════
\begin{exercise}
Prove both directions of Theorem~\ref{thm:equiv-partition} for: congruence mod $n$ on $\Z$; and
$x\sim y\iff |x|=|y|$ on $\R$ (describe the partition geometrically).
\end{exercise}
\begin{exercise}
Carry out the construction of $\Q$ from $\Z\times(\Z\setminus\{0\})$: verify $(a,b)\sim(c,d)\iff ad=bc$
is an equivalence relation, and that $[(a,b)]+[(c,d)]=[(ad+bc,bd)]$ is well defined.
\end{exercise}
\begin{exercise}
Show $f\preceq g\iff f=O(g)$ is a preorder on positive functions $\N\to\R_{>0}$ but not a partial
order; describe the quotient poset (growth rates) and identify its minimum.
\end{exercise}
\begin{exercise}
In $(\Pow(X)\setminus\{\varnothing\},\subseteq)$ with $|X|\ge2$: find minimal and maximal elements;
does a minimum exist? Repeat for $(\N_{\ge2},\mid)$.
\end{exercise}
\begin{exercise}
Prove Proposition~\ref{prop:order-forms} in full (all case analyses), and show lexicographic order on
$\R^2$ is a strict total order.
\end{exercise}
\begin{exercise}
Prove the forward direction of Proposition~\ref{prop:wo-basic} in ZF, and pinpoint exactly where the
converse uses dependent choice; give a poset with no decreasing chain that is not well-ordered.
\end{exercise}
\begin{exercise}
Prove Proposition~\ref{prop:zorn-ac} in full detail (verify $\bigcup\mathcal C$ is a function and a
partial choice function; verify maximality forces full domain).
\end{exercise}
\begin{exercise}
Use Zorn to prove Krull's theorem: every proper ideal of a commutative ring with $1\neq0$ lies in a
maximal ideal (this is S1's engine --- prove it here first).
\end{exercise}
\begin{exercise}
Use Zorn to prove every vector space has a basis (order the linearly independent subsets by
inclusion); identify where the chain upper bound uses unions.
\end{exercise}
\begin{exercise}
Prove Tychonoff's theorem for \emph{finite} products in ZF by induction, and locate the exact step
where the general proof must invoke choice (ultrafilters or subbases).
\end{exercise}
\begin{exercise}
Show AC $\iff$ every surjection $p:E\to B$ has a section $s:B\to E$ ($p\circ s=\id_B$); interpret
sections as ``continuous choices of points of fibers'' and compare with Remark~1.4.
\end{exercise}
\begin{exercise}
Explain why Russell's paradox does not contradict Separation: for a set $A$, show
$\{x\in A:x\notin x\}\notin A$; conclude no universal set exists in ZF.
\end{exercise}

% ══════════════════════════════════════════════
\section*{Where This Session Feeds the Course}
\addcontentsline{toc}{section}{Where This Session Feeds the Course}
% ══════════════════════════════════════════════
\begin{itemize}
  \item \textbf{Depends on:} nothing --- this is session zero.
  \item \textbf{Feeds:} S1 (Krull via Zorn; quotients $R/I$); S2 (tensor as quotient; choosing
    representatives); S3 (algebraic closure via Zorn); S27 (sheafification classes); S34/S16
    (orderings of finite subextensions for inverse limits); S58 (Deligne via Zorn); P1 (cosets as
    partitions, Exercise~12 preview).
  \item \textbf{Unifying thread:} quotient by equivalence $=$ identifying points that must agree;
    order $=$ the scaffolding on which Zorn climbs; Zorn $=$ the engine that guarantees the
    \emph{existence} of the points (maximal ideals, closures, models) the course studies.
\end{itemize}

\subsection*{References for this session}
\begin{enumerate}[label={[\arabic*]}]
  \item Halmos, P.R. \emph{Naive Set Theory} (the classic first read).
  \item Enderton, H.B. \emph{Elements of Set Theory} (proofs of the equivalence cycle).
  \item Jech, T. \emph{The Axiom of Choice} / \emph{Set Theory} (the definitive treatment).
  \item Howard, P. \& Rubin, J. \emph{Consequences of the Axiom of Choice} (the list, with tiers).
  \item Davey, B. \& Priestley, H. \emph{Introduction to Lattices and Order} (order theory).
\end{enumerate}

\vspace{1.5em}
\begin{center}
\rule{0.6\textwidth}{0.4pt}\\[0.5em]
\textit{End of Session P0 --- next: Session P1, Group Theory: Part I}
\end{center}

\end{document}